%% notes-tex/019-simb-multimodal/sections/7-common.tex
%% SOURCE: the ceiling column is assembled from
%% experiments/019-simb-multimodal/results/{expression_ceiling_replicate,
%% morphology_noise_ceiling,pigment_noise_ceiling}.json; the score column from
%% results/round_leaderboards.csv. Both are written by committed scripts.

\section{What the strands have in common\secstatus{todo}}
\label{sec:common}

\subsection{One table, read the same way\secstatus{todo}}

%% GENERATED by experiments/019-simb-multimodal/scripts/build_retrospective_tables.py
%% SOURCE: results/round_leaderboards.csv + the three ceiling JSONs. Do not hand-edit.
\begin{table}[htbp]
\centering\small
\begin{tabular}{lrrrrl}
\toprule
strand & ceiling & best & of ceiling & epochs (peak/last) & runs \\
\midrule
Expression (Kemmeren, Sameith) & 0.775 & \textbf{0.229} & 30\% & 4007/4105 & 163 \\
Expression, masked-label objective & 0.775 & \textbf{0.241} & 31\% & 9188/9997 & 16 \\
Morphology (CalMorph, Ohya) & 0.611 & \textbf{0.082} & 13\% & 27/65 & 23 (1 flat) \\
Expression and morphology, joint & -- & \textbf{0.051} & -- & 28/70 & 32 \\
Betaxanthin (Cachera) & 0.914 & \textbf{0.372} & 41\% & 254/999 & 44 \\
Betaxanthin with metabolome head & 0.914 & \textbf{0.430} & 47\% & 176/496 & 46 \\
Amino-acid pools (Mulleder) & -- & \textbf{0.211} & -- & 402/999 & 31 \\
Beta-carotene (Ozaydin) & 0.544 & \textbf{0.132} & 24\% & 799/999 & 39 \\
\bottomrule
\end{tabular}
\caption[]{Every strand, scored the same way. \emph{best} is \texttt{roll\_max}, the maximum of a centered five-point rolling mean of the validation curve, an upward-biased order statistic whose bias grows with epochs run, which is why the epoch column travels with it. \emph{epochs (peak/last)} gives where the best run peaked and where it stopped: a peak in the last tenth means the run was cut off while still improving. Ceilings are measured by three different estimators because the three measurement types require them; the amino-acid cell is empty because Mulleder has one replicate per strain, and no ceiling is estimable. \emph{flat} counts runs whose validation metric never left zero.}
\label{tab:strand-summary}
\end{table}


\tcfigfit[0.72]{figures/retrospective_achieved_vs_ceiling.pdf}{Achieved score against
measured ceiling, best run per strand. The dashed outline is the ceiling and the gap is the
headroom; amino acid and the expression-morphology joint arm carry no outline because no
single ceiling applies to them. A bare correlation is not comparable across these rows,
which is the reason the panel exists.}{fig:achieved-vs-ceiling}

\tcfigfit[0.72]{figures/retrospective_peak_position.pdf}{Where the best run's peak sits
inside its own budget. Points near the top were still improving when the run stopped, so
their scores are lower bounds set by compute. Points in the lower band peaked early and
spent the rest of the budget overfitting.}{fig:peak-position}

Three things are visible only when the strands sit together. The ceilings differ by more
than a factor of two, so an unqualified $r$ is uninterpretable across rows. The epoch
budgets differ by more than a factor of a hundred, so the score column is partly a record of
how long each strand was allowed to run. And the two strands closest to their ceilings,
betaxanthin and amino acid, are the two with a scalar or low-dimensional target.

\subsection{Two opposite convergence failures, and neither was diagnosed from a curve\secstatus{todo}}

\begin{itemize}
\item \textbf{Expression never peaks.} At every budget tried, up to 10,000 epochs, the
  validation Pearson peaks in the last tenth of the run. Raising the budget from 1,600 to
  10,000 raised the score from $0.204$ to $0.241$ and moved the peak with it.
\item \textbf{The metabolite strands peak early and then overfit.} Amino-acid runs peak at
  epochs 48 to 134 with validation slope turning negative while training slope stays
  positive. The best betaxanthin run peaks at epoch 254 of 999.
\item \textbf{Morphology was never run long enough to tell which it does.} Longest run,
  121 epochs.
\end{itemize}

Same encoder, same optimizer family, opposite behavior. The difference is in the target: a
6,127-dimensional vector of $\log_2$ ratios against a scalar or a 19-vector. Any single
epoch budget applied across strands is wrong for at least two of them, which is what
produced the truncated expression waves and would produce overfit metabolite runs if the
budget were set from expression.

\subsection{Loss and metric point in different directions on expression alone\secstatus{todo}}

Validation loss bottoms 40 to 250 epochs before validation Pearson peaks in wave 4b, and
$\approx 1{,}200$ epochs before it in the long runs, while the prediction
standard-deviation ratio climbs from $0.001$ to $0.52$. The 010 fitness and interaction task
does not do this: there both losses fall monotonically and the validation-loss minimum is at
the last epoch. The consequence already cost the project every checkpoint the expression
strand saved before best-by-metric checkpointing landed.

\notesec{What this suggests, unverified} A quantile loss on a $\log_2$ ratio can be
minimized by a narrow, well-calibrated predictive distribution that gets the ordering
wrong, and Pearson is invariant to the global rescale that the loss is not. Reporting
\texttt{nmse} alongside, defined so that $\texttt{nmse} = 1$ is exactly ``predict each
gene's mean'', separates right-ordering from right-magnitude. That instrument exists; the
experiment that would test whether a different loss closes the gap does not.

\subsection{Configuration choice moves results more than method choice\secstatus{todo}}

On betaxanthin the ten grid cells span Spearman $0.013$ to $0.158$ and AUC $0.406$ to
$0.695$, against a between-method gap of $0.0014$ on AUC. Four low points are the
$\mathrm{lr} = 10^{-3}$ cells, which collapsed to a constant predictor. On the amino-acid
strand the best run is $0.211$ and the median is $0.0498$.

Any single-number claim about a method is dominated by which cell is quoted, which is why
the cell must be fixed by validation in advance, and why the leaderboard records the median
alongside the maximum. It also means an arm difference smaller than the cell spread is not
measurable without pinning the cell and replicating seeds inside it.

\subsection{Cross-modality overlap is real and small\secstatus{todo}}

Before spending compute on joint proteome and expression training, the linear overlap
between the two was measured on the 1,350 strains they share. Per-gene correlation across
strains has median $r = 0.08$, with $1.7\,\%$ of genes above $0.3$ and $8.9\,\%$ negative.
A ridge map recovers held-out $R^2$ of $0.035$ in one direction and $0.024$ in the other,
above a trivial baseline of $-0.004$ but far from redundancy.
\src{experiments/019-simb-multimodal/results/proteome_expression_eda.json}

\tcfigfit{figures/proteome_expression_per_gene_corr_hist_2026-07-22-17-15-26.pdf}{Per-gene
correlation across strains between the Messner proteome and the Kemmeren expression
compendium, over the 1,832 genes and 1,350 strains they share. A narrow lump barely off
zero. The few high-$r$ genes are dominated by the strain in which that gene is itself
deleted, which drives both its transcript and its protein toward zero.}{fig:prot-expr}

The two modalities were measured in different media, synthetic minimal for the proteome and
synthetic complete for the expression, which attenuates every number above. The overlap is
therefore certified as non-redundant and not certified as cleanly complementary. That
distinction is what bounds how hard a ``the proteome helps expression'' claim can be pushed
until a same-media pair exists.

\subsection{The graph prior has never been checked\secstatus{todo}}

The nine interaction graphs enter the model as a hard additive attention mask, so composing
$L$ layers makes the influence of gene $X$ on gene $Y$ approximately the $L$-step
random-walk reachability of $Y$ from $X$ on that graph. That is a falsifiable assumption
about the data: \emph{deleting $X$ perturbs the genes near $X$ on graph $k$}. If it is
false, pulling attention toward the graph pushes the model toward the wrong prior and no
setting of the regularization weight rescues it.

The probe that would test it is specified and unbuilt. It costs minutes of CPU: for each
deleted gene, the across-reporter rank correlation between graph proximity and response
magnitude, reported per graph as an AUC against degree-preserving rewiring and a matched
random graph, where $0.5$ means the graph says nothing.
\src{notes/experiments.019-simb-multimodal.graph-prior-probe.md}

Nine graphs consume nine attention heads. If three carry signal, the other six are
constrained toward noise and are worse than free heads. This is the cheapest unrun
experiment in the project.

\subsection{The perturbation axis is the binding resource\secstatus{todo}}

\Cref{tab:label-accounting} makes the point for expression: 9.2 million scalar labels
across 1,484 perturbations, each gene deleted exactly once. Morphology is the same shape,
1.3 million labels across 4,718 perturbations, each once. The trigenic interaction data is
the opposite shape, 91 thousand labels across 91 thousand records in which a given gene
recurs with many partners.

The strands that work in this project are the ones where the supervision constrains a gene
repeatedly. \textbf{Hypothesis, untested:} the reason multi-task training is attractive here
is not that the labels are informative about each other, but that a second phenotype gives
a second, differently-shaped constraint on the same perturbation. That predicts the gain
should be larger where the two phenotypes share strains and smaller where they merely share
genes, which is a testable statement and has not been tested.
